\section{Why Now}

Several trends make execution optimization timely.

First, provider fragmentation is increasing. Organizations use multiple frontier providers, local models, specialized models, and gateways. Provider-neutral execution policy becomes more valuable as the provider landscape becomes heterogeneous.

Second, context is larger but not free. Large context windows reduce some pressure but increase hidden waste. Prompt caching makes context layout itself an execution optimization problem \citep{openai_prompt_caching,anthropic_prompt_caching}.

Third, agents are becoming long-running. Durable execution, state persistence, and human-in-the-loop workflows are increasingly part of agent infrastructure \citep{langgraph_overview,temporal_agents,openai_agents_running}.

Fourth, tool boundaries are becoming protocolized. MCP turns tools into structured, discoverable interfaces rather than one-off function blobs \citep{mcp_tools}.

Fifth, observability and evaluation are maturing. Traces, cost metrics, latency metrics, prompt management, datasets, and evaluator outputs are becoming production expectations \citep{langfuse_docs,otel_genai}.

Together, these trends create the preconditions for an optimizer. Once execution becomes structured, measured, and multi-provider, optimization becomes possible.
